\section{Limitations}
\label{sec6.2: limitations}

Despite the empirical results, this framework has several methodological and biological limitations:

\begin{itemize}
    \item \textbf{Small Sample Host Instability:} Training sample sizes vary widely across hosts (ranging from
          $>1\,400$ phages for \emph{Mycobacterium} down to fewer than $10$ for rare genera such as \emph{Sulfolobus}
          and \emph{Microcystis}). Abundant hosts calibrate reliably near the $90\%$ target, but low sample hosts
          show significantly higher empirical variance, since the quantile estimates on small calibration sets ($n_h < 10$)
          struggle to capture the tails of the true infector score distribution. If we restrict our evaluation to hosts
          with $\geq 100$ training phages, we can recover coverage for the Strip method (\ref{sec5.3.6: host ablations}), 
          but this ablation also reduced the bin count from the validation-tuned $N_B = 10$ to $N_B = 5$ in the same step, 
          so we cannot attribute the improvement just to sample size.
         
    \item \textbf{High Dimensional Sensitivity:} Collapsed 1D achieves compact prediction sets by 
          compressing score vectors into a scalar mean, making it unable to distinguish a phage whose evidence concentrates in 
          one informative dimension from one whose evidence is spread thinly across many uninformative ones (Fig \ref{fig: 
          staph 1 good 1 bad}). Conversely, the Radial and Strip envelopes preserve this directional structure but suffer as 
          $K_h$ grows toward $40$: sparse angular or bin-wise support forces frequent fallback to the conservative marginal 
          quantile, which is the direct cause of their larger average prediction sets relative to Collapsed 1D.

    \item \textbf{Structural Missingness and Marginal Fallbacks:} Phages frequently lack entire functional protein classes, 
           producing $30\%\text{--}36\%$ missing values per score vector. These missing coordinates force the multi dimensional 
           envelopes to fall back on marginal quantiles rather than the tighter conditional bounds, which widens prediction 
           boundaries and increases co-prediction among closely related host genera (e.g.\ within \emph{Enterobacteriaceae}) 
           that share similar missingness patterns.

    \item \textbf{Large Absolute Set Sizes:} Even where coverage is valid, the prediction sets of Radial and Strip are on average 
           os size $15+$ out of $54$ candidate hots. A list of that size is of limited direct use as a clinical or laboratory 
           shortlist, even though it is a large efficiency gain over the unrestricted, non Gram filtered setting (avg.\ size 
           $\approx 39$, Table \ref{tab: gram filtering effect}). 

    \item \textbf{Error Propagation:} Because host level prediction depends on Gram type filtering, the theoretical miscoverage 
           is added ($\alpha + \max_h \alpha_h$, Proposition \ref{prop:cascade}). Although the empirical Gram type miscoverage 
           was near zero in our evaluation, any systematic bias in the first stage classification, (for example, on phages with 
           ambiguous raw scores near the $0.25\text{--} 0.75$ zone) prunes the true host before the host level envelope is even 
           evaluated, and this cannot be corrected.

\end{itemize}

